\section{先认清要估计的数}
\subsection{用素数幂算一个偏差}
取一个尺度 $x$。对每个不超过 $x$ 的素数幂 $p^r$，按下面的规则记一项：
\[
       \left(\frac{x}{p^r}-1\right)\log p.
\]
素数幂越靠近 $x$，这项的权重越小；恰等于 $x$ 时贡献为零。比如 $x=6$，只用 $2,3,4,5$ 四个素数幂，算出来是
\[
 S(6)=2\log2+\log3+\tfrac12\log2+\tfrac15\log5.
\]
$4=2^2$ 那一项仍乘 $\log 2$，不是 $\log4$。把这条规则写成通式，令 $\Lambda(p^r)=\log p$，其余 $\Lambda(n)=0$，则
\[
 \begin{aligned}
 S(x)&=\sum_{1\le n\le x}\Lambda(n)(x/n-1),\\
 G(x)&=S(x)-\bigl[x\log x-(1+\gamma_E)x\bigr].
 \end{aligned}
\]
$\gamma_E$ 是欧拉常数。方括号是用作比较的主项；$G(x)$ 才是我们说的“偏差”，即实际算出的数减去主项。选择这种权重和主项，与 $\zeta$ 函数的积分公式有关，见本节末。\src{Development/AutomorphicArithmetic/DedekindLFunctions/CharFamily.lean}{L10-L22}{权重与偏差定义}。

本章不是直接证明对每个 $x$ 都有某个 $|G(x)|$ 上界。实际论证还要引入加权的偏差组合，并在合适的长区间内控制它的单侧上界。第 5 节写出区间条件；从它到零点还要用解析判据。先记住这里的分工：$G(x)$ 是解释最终对象的起点，后文的局部收益不是把 $G(x)$ 换个名字。

\subsection{两条证据在哪里用到一起}
先用整数矩阵把接口说清。假如一个矩阵的两列是 $(2,0)$、$(4,0)$，每个坐标都是整数，但两列只有一个独立方向。另一个矩阵的两列是 $(1/2,0)$、$(0,1/2)$，方向有两个，却不能把所有整数坐标送成整数坐标。前一个例子表明\textbf{整性不推出满秩}，后一个表明\textbf{满秩不推出整性}。

在真正的证明里，也把同一个映射写成矩阵。一条路证明：先定下同一个检验格，再按每个有限处的整除预算调整输入格，证明映射把调整后的输入送进指定的整对偶格。另一条路证明：用复矩阵控制这些映射的大小；删去一些输入方向后，把各个映射剩余的像合起来，仍能给出独立方向数的下界。“像”就是所有输出组成的空间，“秩”是其中独立方向的个数。

这两条路在第 3 节会合。那里要问：删掉一部分输入方向之后，剩下的像在\textbf{同一个输出格}里还占多大？要回答它，既要知道像有多少独立方向，又要知道这些方向确实来自指定格里的输出。换成另一个矩阵，即使各自的统计量都好看，也不能一起用。到了第 5 节，两个不等式从上下夹住同一个量：按各输入格的秩加权后的斜率。比较上下界时，同一个输出格的平均斜率消去，留下要控制的算术差。这就是矩阵那一支究竟补了什么。

\subsection{进入正文前，只认三个对象}
最简单的格是平面上的整数点 $\mathbb Z^2$。这里的格也规定哪些系数可以取，只是整数换成了数域的整数环。下面用到的三个对象是输入格、检验格和映射。做比较时固定一组数据。用 $u$ 编号不同的算术链，源码叫它“词”；这里先把它当作输入格的编号。

写 $R$ 为公共检验格，$C=R^\vee$ 为整对偶格。对偶这里的含义很具体：线性泛函取 $R$ 的任意元素，结果都在整数环。源侧从原格 $N_u^0$ 变成理想平移后的 $N_u$。只有在平移后才能写出本章要用的包含
\[
       B_u(N_u)\subseteq C.
\]
记 $h_u$ 为同一个 $B_u$ 的高度上界。证明不会从整性直接推出高度，也不会从高度推回整性。本章使用 $B_u$ 同时表示双线性配对与把第二输入留空后的线性映射，亦即 $\phi\mapsto (F\mapsto B_u(\phi,F))$。


\subsection{补读：为什么选择这一种权重}
熟悉复分析的读者可以用下式核对权重的来源；跳过本段不影响第 2 节。令 $W(t)=(t-1)_+$，在 $\Re z>1$ 时
\[
 \int_1^\infty(t-1)t^{-z-1}dt
  =\frac1{z-1}-\frac1z=\frac1{z(z-1)}.
\]
取 $c>1$。反演这个 Mellin 变换，再用等式
\[
 -\zeta'(z)/\zeta(z)=\sum_n\Lambda(n)n^{-z},
\]
便得到
\[
 S(x)=\frac{1}{2\pi i}\int_{c-i\infty}^{c+i\infty}
     -\frac{\zeta'(z)}{\zeta(z)}
       \frac{x^z}{z(z-1)}\,dz .
\]
这里可交换绝对收敛的级数与积分，因为在直线 $\Re z=c$ 上，两项条件
$\sum_n\Lambda(n)n^{-c}<\infty$ 与
$\int_\Rr |(c+it)(c-1+it)|^{-1}dt<\infty$ 同时成立。若 $x$ 恰是整数，边界权重为零，无需决定是否给最后一项半权。

这给出了一点具体的预期：$z=1$ 处，对数导数的极点与分母的极点重叠，二阶极点带出 $x\log x$ 型主项；若考察非平凡零点 $\rho$，其指数尺度是 $x^{\Re\rho}$。但这点直觉\emph{不是}本证明已用的最终零点推理。上界只有一侧，不同零点的项还会互相抵消；把每项的尺度和全体零点的结论直接连起来，正是需要避免的跳步。后面以计划好的窗口和解析准则来处理，而不是在本章就把这条积分线移到任意指定位置。

